\chapter{Resultados y Análisis}

\section{Primera parte: Viga suspendida del orificio <<A>>}
Al suspender la viga desde su centro de gravedad, el torque ejercido por la propia viga es cero ($\tau_{\text{viga}} = 0$).

\begin{table}[H]
\centering
\begin{tabular}{lccc}
\toprule
\textbf{Parámetros} & \textbf{Caso <<d>>} & \textbf{Caso <<e>>} & \textbf{Caso <<f>>} \\
\midrule
Masa 1 (kg) & 0.1 & 0.1 & 0.1 \\
Radio 1 (m) & 0.15 & 0.3 & 0.3 \\
\textbf{Torque 1}* (Nm) & 0.147 & 0.294 & 0.294 \\
\midrule
Masa 2 (kg) & 0.203 & 0.203 & 0.203 \\
Radio 2 (m) & 0.3 & 0.15 & 0.15 \\
\textbf{Torque 2}* (Nm) & 0.59682 & 0.29841 & 0.29841 \\
\midrule
Masa 3 (kg) & 0.303 & 0.303 & 0.564 \\
Radio 3 (m) & 0.24 & 0.2 & 0.1 \\
\textbf{Torque 3}* (Nm) & 0.0712656 & 0.59388 & 0.55272 \\
\midrule
Torque viga (Nm) & 0 & 0 & 0 \\
\textbf{Sumatoria de torques} (Nm) & \textbf{0.031164} & \textbf{-0.00147} & \textbf{0.03969} \\
\bottomrule
\end{tabular}
\caption{Datos correspondientes a la primera parte del experimento.}
\end{table}

\newpage

\section{Segunda parte: Viga suspendida del orificio <<B>>}
Al suspender la viga desde el orificio <<B>>, el centro de masa de la viga ejerce un torque adicional de 0.154 Newtons por metro que interactúa con las demás masas para mantener el sistema horizontal.

\begin{table}[H]
\centering
\begin{tabular}{lccc}
\toprule
\textbf{Parámetros} & \textbf{Caso <<g>>} & \textbf{Caso <<h>>} & \textbf{Caso <<i>>} \\
\midrule
Masa 1 (kg) & 0.048 & 0.048 & 0.048 \\
Radio 1 (m) & 0.05 & 0.1 & 0.2 \\
\textbf{Torque 1}* (Nm) & 0.02352 & 0.04704 & 0.09408 \\
\midrule
Masa 2 (kg) & 0.101 & 0.101 & 0.101 \\
Radio 2 (m) & 0.15 & 0.05 & 0.05 \\
\textbf{Torque 2}* (Nm) & 0.14847 & 0.04949 & 0.04949 \\
\midrule
Masa 3 (kg) & 0.82 & 0.82 & 0.82 \\
Radio 3 (m) & 0.2 & 0.2 & 0.2 \\
\textbf{Torque 3}* (Nm) & 1.6072 & 1.6072 & 1.6072 \\
\midrule
Torque viga* (Nm) & 1.5092 & 1.5092 & 1.5092 \\
\textbf{Sumatoria de torques} (Nm) & \textbf{-0.07399} & \textbf{0.00147} & \textbf{-0.04557} \\
\bottomrule
\end{tabular}
\caption{Datos obtenidos durante la segunda parte del experimento.}
\end{table}

